% =============================================================================
% cap03.tex — REVISAO BIBLIOGRAFICA
% =============================================================================

\chapter{REVISÃO BIBLIOGRÁFICA}\label{cap:revisao}

Esta revisão percorre os fundamentos que sustentam as verificações executadas
no trabalho. Parte da constituição do material, passa pelas propriedades de
cálculo e pelos coeficientes que as modificam, e chega às prescrições
específicas aplicáveis a vigas de seção variável, à estabilidade, ao
contraventamento e às ligações.

\section{A MADEIRA LAMELADA COLADA COMO PRODUTO ESTRUTURAL}\label{sec:mlc-produto}

A MLC é obtida pela colagem de lamelas de madeira serrada, classificadas
previamente quanto à resistência, com as fibras paralelas ao eixo do elemento.
A EN 14080 \cite{EN14080:2013} fixa os requisitos de produção e as propriedades
características do produto acabado. As lamelas são emendadas topo a topo por
entalhes múltiplos, o que permite comprimentos muito superiores ao da madeira
serrada disponível, e são coladas em prensa com adesivos estruturais
classificados pela EN 301 \cite{EN301:2017}.

Três consequências desse processo interessam ao projeto estrutural. A primeira
é a liberdade geométrica: a espessura reduzida das lamelas permite curvá-las
antes da cura e produzir elementos de altura variável ao longo do vão,
ajustando a seção ao diagrama de momentos. A segunda é a homogeneização: os
defeitos naturais da madeira, sobretudo os nós, são distribuídos
aleatoriamente entre as lamelas em vez de se concentrarem em uma única seção
crítica, o que reduz a dispersão da resistência e eleva os valores
característicos. A terceira é a estabilidade dimensional, decorrente da
secagem controlada das lamelas antes da colagem.

A EN 14080 \cite{EN14080:2013} organiza o produto em classes designadas pela
letra GL seguida da resistência característica à flexão em megapascal, com o
sufixo \textit{h} para composição homogênea, em que todas as lamelas pertencem
à mesma classe, e \textit{c} para composição combinada, em que lamelas de
classe superior ocupam as zonas mais solicitadas da seção.

%% TODO: referência — acrescentar literatura sobre o processo produtivo da MLC
%% e sobre a eficiência das emendas dentadas (ex.: Porteous \& Kermani;
%% Herzog et al.; Calil Junior). Chaves a criar em referencias.bib.

\section{PROPRIEDADES DE CÁLCULO}\label{sec:propriedades-calculo}

O Eurocódigo 5 obtém o valor de cálculo de uma propriedade de resistência a
partir do valor característico por meio da expressão

\begin{align}
    X_d = k_{mod}\,\frac{X_k}{\gamma_M} \label{eq:valor-calculo}
\end{align}

\noindent
em que $\gamma_M$ é o coeficiente parcial do material, igual a 1,25 para
elementos de MLC e a 1,30 para ligações, e $k_{mod}$ é o coeficiente de
modificação.

O coeficiente $k_{mod}$ concentra duas influências que a madeira sofre e que os
materiais isotrópicos usuais não apresentam na mesma intensidade: o teor de
umidade de equilíbrio e a duração do carregamento. Quanto mais úmido o ambiente
e quanto mais prolongada a ação, menor a resistência mobilizável.

A umidade é tratada por classes de serviço. A classe 1 corresponde a ambientes
internos aquecidos, com umidade de equilíbrio da ordem de 12\%; a classe 2, a
ambientes cobertos porém não aquecidos, com umidade até cerca de 20\%; e a
classe 3, a condições mais severas que as anteriores, típicas de peças expostas
à intempérie ou em edificações abertas. A duração é tratada por cinco
categorias, que vão da ação permanente, superior a dez anos, à ação
instantânea. A \autoref{tab:kmod-kdef} apresenta os valores aplicáveis à MLC.
Quando a combinação reúne ações de durações distintas, adota-se o $k_{mod}$ da
ação de menor duração entre as presentes.

\tabela
{Coeficientes de modificação e de fluência para madeira lamelada colada}
{\footnotesize
\begin{tabular}{ccccccc}
\toprule
    \textbf{Classe de} & \multicolumn{5}{c}{\textbf{$k_{mod}$ por classe de duração da ação}} & \textbf{$k_{def}$} \\
\cmidrule(lr){2-6}
    \textbf{serviço} & Permanente & Longa & Média & Curta & Instantânea & \\
\midrule
    1 & 0,60 & 0,70 & 0,80 & 0,90 & 1,10 & 0,60 \\
    2 & 0,60 & 0,70 & 0,80 & 0,90 & 1,10 & 0,80 \\
    3 & 0,50 & 0,55 & 0,65 & 0,70 & 0,90 & 2,00 \\
\bottomrule
\end{tabular}}
{Adaptado da EN 1995-1-1 \cite{EN1995-1-1:2004}, tabelas 3.1 e 3.2}
{kmod-kdef}
{}
{}

O coeficiente $k_{def}$ quantifica a fluência e é empregado nas verificações de
estado-limite de serviço para converter a flecha instantânea na flecha final. A
progressão de 0,60 na classe 1 até 2,00 na classe 3 mostra que a deformação
diferida de uma peça exposta pode triplicar a parcela instantânea, o que faz da
classe de serviço uma decisão de projeto com consequência direta sobre a altura
da seção.

\section{VIGAS DE SEÇÃO VARIÁVEL}\label{sec:vigas-secao-variavel}

A viga de duplo-bisel tem altura crescente do apoio até a cumeeira, situada a
meio vão, acompanhando o crescimento do momento fletor. Essa eficiência tem
contrapartidas que o Eurocódigo 5 trata em prescrições próprias.

A primeira decorre da borda inclinada. Em uma viga de altura constante, as
tensões de flexão são paralelas às faces. Quando uma face é inclinada, a tensão
na borda deixa de ser uniaxial e passa a combinar componentes paralela e
perpendicular às fibras, além de cisalhamento. O item 6.4.2 da
EN 1995-1-1 \cite{EN1995-1-1:2004} trata essa situação por meio do coeficiente
$k_{m,\alpha}$, que reduz a resistência à flexão na borda inclinada. A redução é
assimétrica: é mais severa quando a borda inclinada está tracionada do que
quando está comprimida, porque a resistência à tração perpendicular às fibras é
muito inferior à resistência à compressão na mesma direção.

A segunda decorre da zona de cumeeira. Na região do vértice, a mudança de
direção das tensões de flexão gera tensões de tração perpendicular às fibras,
que constituem a fragilidade característica da madeira. O item 6.4.3 da mesma
norma verifica essas tensões introduzindo o coeficiente $k_{dis}$, que responde
pela distribuição não uniforme, e um fator de volume, que incorpora o efeito de
escala: quanto maior o volume solicitado, maior a probabilidade de conter um
defeito capaz de iniciar a ruptura.

\section{ESTABILIDADE}\label{sec:estabilidade-revisao}

A esbeltez das vigas de MLC torna a estabilidade lateral, e não a resistência
da seção, o critério que frequentemente governa. O item 6.3.3 da
EN 1995-1-1 \cite{EN1995-1-1:2004} trata a flambagem lateral por torção
mediante o coeficiente $k_{crit}$, que reduz a resistência à flexão em função da
esbeltez relativa. Essa esbeltez depende do comprimento efetivo entre pontos de
contenção lateral, o que estabelece a ligação direta entre a verificação da
viga e o sistema de contraventamento: terças efetivamente vinculadas a um
diafragma contraventado reduzem o comprimento efetivo ao espaçamento entre
elas, enquanto terças não contraventadas deixam a viga com o vão inteiro como
comprimento efetivo.

Para os pilares, sujeitos a esforço normal e momento simultâneos, aplicam-se as
expressões de flexo-compressão do item 6.3.2, que combinam as taxas de
utilização por meio dos coeficientes $k_c$ relativos a cada eixo principal.

Em nível de sistema, a estabilidade global é avaliada pelo fator $\alpha_{cr}$,
razão entre a carga crítica elástica e a carga de cálculo. Valores de
$\alpha_{cr}$ superiores ao limite normativo permitem dispensar a análise de
segunda ordem e tratar a estrutura como indeslocável, desde que consideradas as
imperfeições geométricas por meio de forças horizontais equivalentes.

\section{CONTRAVENTAMENTO}\label{sec:contraventamento-revisao}

O item 9.2.5 da EN 1995-1-1 \cite{EN1995-1-1:2004} estabelece as forças que um
sistema de contraventamento deve resistir para estabilizar um conjunto de
elementos comprimidos ou fletidos. A formulação parte da imperfeição inicial
dos elementos estabilizados e conduz a uma carga distribuída sobre o sistema
contraventante, modulada pelo coeficiente

\begin{align}
    k_{\ell} = \min\left(1;\ \sqrt{\frac{15}{\ell}}\right) \label{eq:k-ell}
\end{align}

\noindent
em que $\ell$ é o vão do sistema de contraventamento em metros. O sistema é
paralelo aos elementos que estabiliza e se apoia nos mesmos pontos que eles, de
modo que, no caso de vigas, $\ell$ coincide com o vão das próprias vigas, e não
com o comprimento da edificação nem com a distância entre os sistemas. O número
de sistemas contraventados não altera $\ell$: altera apenas o número $n$ de
elementos que cada sistema estabiliza, e mais sistemas significam menor força
em cada um.

\section{LIGAÇÕES}\label{sec:ligacoes-revisao}

O dimensionamento de ligações por pinos metálicos apoia-se na teoria de
Johansen, incorporada ao item 8.2 da EN 1995-1-1 \cite{EN1995-1-1:2004}. O
método examina os mecanismos cinematicamente admissíveis de uma ligação, que
combinam o esmagamento da madeira sob o pino e a formação de rótulas plásticas
no próprio pino, e adota como resistência o menor valor entre eles. Os modos
que envolvem plastificação do pino são dúcteis e permitem redistribuição de
esforços.

A esses mecanismos somam-se modos de ruptura frágeis, que não admitem
redistribuição e cuja ocorrência é súbita. O item 8.1.4 trata do fendilhamento
provocado por componentes de força perpendiculares às fibras, e o anexo A trata
da ruptura de bloco, em que um volume de madeira delimitado pelo grupo de pinos
se destaca da peça. O item 8.1.2 introduz ainda o número efetivo de pinos
$n_{ef}$, inferior ao número real quando os pinos se alinham paralelamente às
fibras, reconhecendo que a distribuição de carga entre eles não é uniforme. A
\autoref{tab:modos-ruptura} resume essa classificação, que orienta a adoção de
taxas de utilização máximas distintas para cada natureza, conforme exposto no
\autoref{cap:detalhamento}.

\tabela
{Classificação dos modos de ruptura em ligações de madeira}
{\footnotesize
\begin{tabular}{lll}
\toprule
    \textbf{Natureza} & \textbf{Modo de ruptura} & \textbf{Prescrição} \\
\midrule
    Dúctil  & Plastificação do pino com esmagamento da madeira & Item 8.2 \\
    Dúctil  & Escoamento da chapa metálica                     & EN 1993-1-8 \\
    Frágil  & Esmagamento localizado sem plastificação do pino  & Item 8.2 \\
    Frágil  & Fendilhamento por tração perpendicular às fibras  & Item 8.1.4 \\
    Frágil  & Ruptura de bloco no grupo de pinos                & Anexo A \\
\bottomrule
\end{tabular}}
{Elaborado pelo autor (2026) com base na EN 1995-1-1 \cite{EN1995-1-1:2004}}
{modos-ruptura}
{}
{}

\section{AÇÕES E COMBINAÇÕES}\label{sec:acoes-revisao}

As ações permanentes e as sobrecargas de utilização seguem a
EN 1991-1-1 \cite{EN1991-1-1:2002}, e as ações durante a execução, relevantes
para a fase de içamento, seguem a EN 1991-1-6 \cite{EN1991-1-6:2005}. As
combinações de ações e os coeficientes parciais e de combinação decorrem da
EN 1990 \cite{EN1990:2002}.

A ação do vento merece tratamento à parte. Embora o Eurocódigo 1 disponha de
parte específica sobre o assunto, a EN 1991-1-4 \cite{EN1991-1-4:2005}, seus
parâmetros descrevem o regime de ventos europeu e dependem de mapas de
velocidade básica inexistentes para o território brasileiro. Adota-se, por isso,
a ABNT NBR 6123 \cite{NBR6123:2023}, que determina a velocidade característica
a partir da velocidade básica do sítio e dos fatores topográfico, de rugosidade
e dimensões da edificação, e estatístico. Sobre a pressão dinâmica resultante
incidem os coeficientes de pressão externa e interna tabelados conforme a forma
da edificação, sendo que edificações abertas ou com dominante face permeável
admitem coeficientes internos que agravam substancialmente a sucção na
cobertura.
